\documentclass[preview]{standalone}
\usepackage{amsmath}
\usepackage{amssymb}
\usepackage{stellar}
\usepackage{definitions}
\usepackage{bettelini}
\begin{document}
\id{diffeq-theory-exercises}
\genpage
\section{Exercises}

\begin{snippetexercise}{diffeq-nonlipschitz-unique-exercise}{Lipschitz continuity is sufficient, not necessary}
    Find a continuous autonomous scalar field that is not locally Lipschitz
    at zero but for which the initial value problem at zero is unique.
\end{snippetexercise}

\begin{snippetsolution}{diffeq-nonlipschitz-unique-solution}{A logarithmic example}
    Set \(g(0)=0\) and \(g(y)=y\ln|y|\) for \(y\neq0\).
    It is continuous, but \(|g(y)-g(0)|/|y|=|\ln|y||\) is unbounded near zero.
    On a nonzero branch, \(v=\ln|y|\) satisfies \(v'=v\), so
    \[
        y(t)=\operatorname{sign}(y_0)
        \exp\bigl(\ln|y_0|\,e^{t-t_0}\bigr).
    \]
    Such a branch cannot reach zero at any finite time.
    If a solution through zero were nonzero somewhere, take the component
    of its nonzero set containing that point. Its finite boundary toward
    the initial time would force this formula to tend to zero, impossible.
    Thus the only solution through zero is the constant zero solution.
\end{snippetsolution}

\begin{snippetexercise}{diffeq-initial-stability-exercise}{Perturbing only the initial value}
    For a continuous field \(f(t,y)\) with uniform state Lipschitz constant \(L\),
    prove that solutions through \(u(t_0)=u_0\), \(v(t_0)=v_0\) satisfy
    \(\|u(t)-v(t)\|\leq e^{L|t-t_0|}\|u_0-v_0\|\).
\end{snippetexercise}

\begin{snippetsolution}{diffeq-initial-stability-solution}{Integral comparison}
    Subtract the Volterra equations. For \(t\geq t_0\),
    \[
        \|u(t)-v(t)\|\leq\|u_0-v_0\|
           +L\int_{t_0}^t\|u(s)-v(s)\|\,ds.
    \]
    Integral Gronwall proves the estimate. Reverse time for \(t\leq t_0\).
    With local Lipschitz constants, apply the same argument inside a compact
    tube around the limiting solution, as in the continuous-dependence theorem.
\end{snippetsolution}
\end{document}
